\documentclass[../main.tex]{subfiles}
\begin{document}

\chapter{仮想記憶}
\lessoninfo{
  video={Lesson 13，動画 1--27},
  textbook={H\&P \cite{hennessy2017} §B.4--B.5},
  keywords={仮想記憶，物理メモリ，ページ，ページフレーム，ページテーブル，ページフォールト，仮想ページ番号，ページオフセット，多段ページテーブル，ページサイズ，内部断片化，TLB，TLB ミス，ページテーブルウォーク},
}

第12章では，プロセッサが生成するアドレスをすべて主記憶上の位置として
扱った．実際のシステムでは，プログラムが使うアドレスは\emph{仮想}である．
すなわち，どのプログラムも自分専用の大きなアドレス空間を見ており，
オペレーティングシステムとプロセッサが協力して，そのうちプログラムが実際に
使う部分を搭載されたメモリに対応づける．この章では，なぜ仮想記憶が必要か，
ページテーブルがどのように仮想アドレスを物理アドレスに変換するか，それらの
テーブルがどれほど大きく，多段ページテーブルがどのようにそれを小さく保つか，
そして，プロセッサが最近の変換を TLB に保持しなければ，変換がすべての
メモリアクセスを遅くしてしまうのはなぜかを述べる．第14章は，この TLB を
第12章のキャッシュと組み合わせる．

\section{なぜ仮想記憶が必要か}
\label{sec:l13-why}

ハードウェアとソフトウェアでは，メモリの見え方が異なる．\source{Lesson 13，
動画 1--3；H\&P §B.4；\cite{kilburn1962}．}ハードウェアにとって，メモリとは
マシンに搭載されているもの，すなわちメモリモジュールが提供するバイト数だけの
配列である．プログラムにとって，メモリとは 32 ビットプロセッサなら $2^{32}$
バイト，64 ビットプロセッサなら $2^{64}$ バイトのアドレス空間であり，その
すべてがそのプログラムのものである．コード，データ，スタックは，搭載された
メモリの量にも，同時に走る他のプログラムにも関係なく，プログラムとコンパイラ
が選んだアドレスに置かれる．

この2つの見方は，2つの点で食い違う．
\begin{itemize}
  \item 1つのプログラムのアドレス空間は，搭載されたメモリよりはるかに
    大きい．32 ビットのアドレス空間は 4~GB であり，マシンがそれだけの
    メモリをもつとは限らない．64 ビットのアドレス空間は 16~EB
    （1~EB $= 2^{60}$ バイト）で，どのマシンがもつ量よりも何桁も大きい．\caution{64 ビット
    すべてを変換するプロセッサは存在しない．x86-64 は，ビット 63--48 が
    ビット 47 の繰り返しである\emph{カノニカル}アドレスを要求するので，
    使える空間は 16~EB ではなく 256~TB である．}
  \item 多数のプログラムが同時に走り，そのそれぞれが，アドレス空間全体を
    所有しているかのように同じアドレスを使う．
\end{itemize}

\begin{definition}[仮想記憶]
  すべてのプログラムに，大きな自分専用のメモリという見え方を与え，そのうち
  プログラムが実際に使う部分をマシンに搭載されたメモリに対応づける機構を
  \term[かそうきおく]{仮想記憶}[virtual memory]という．
\end{definition}

各プログラムの各部分をどこに置くかはオペレーティングシステムが決め，
プロセッサはその決定をすべてのメモリアクセスのすべてのアドレスに適用する．
これにより，プログラムはメモリの大きさも，一緒に走ることになる他の
プログラムも知らずに書ける．また，オペレーティングシステムが取り計らわない
限り，あるプログラムが別のプログラムのメモリにアクセスすることはできず，
さらに，プログラム全体では搭載量を超えるメモリを使うことができ，あふれた分は
ディスクに置かれる（\cref{sec:l13-page-faults}）．\sidenote{この着想は
それほど古い．1962 年のマンチェスター Atlas は，16K ワードのコア記憶と
96K ワードの磁気ドラムを 512 ワードのページからなる1つのアドレス空間として
見せ，32 個の連想的な\emph{ページアドレスレジスタ}でアドレスを変換した．
\cref{sec:l13-tlb} の TLB の先祖である．}

\section{物理メモリと仮想アドレス空間}
\label{sec:l13-views}

2つの見方は，2種類のアドレスに対応する．\source{Lesson 13，動画
4--5；H\&P §B.4．}

\begin{definition}[物理メモリ]
  マシンに実際に搭載されているメモリを，そのマシンの
  \term[ぶつりめもり]{物理メモリ}[physical memory]という．その各バイトには
  0 から物理メモリの大きさより 1 小さい値まで連番が付けられ，バイトを識別する
  この番号を，そのバイトの\term[ぶつりあどれす]{物理アドレス}[physical address]%
  という．
\end{definition}

物理メモリの大きさを決めるのは搭載されたモジュールであって，プロセッサの
アドレスの幅ではない．32 ビットのマシンは 4~GB をアドレス指定できるが，
搭載されているのは 2~GB かもしれない．64 ビットのマシンはおそらく数十
ギガバイトをもつが，これはアドレス指定可能な $2^{64}$ バイトのごく一部で
ある．この章の例では，物理アドレスのビット数は搭載されたメモリが要求する
ちょうどの数とする．2~GB なら 31 ビットである．

\begin{definition}[仮想アドレス空間]
  プログラムの\term[かそうあどれすくうかん]{仮想アドレス空間}[virtual address space]%
  とは，物理メモリとは無関係に，そのプログラムが使えるアドレスの全範囲，
  すなわち $n$ ビットプロセッサでは 0 から $2^{n}-1$ までである．その中の
  アドレスを\term[かそうあどれす]{仮想アドレス}[virtual address]という．
\end{definition}

プログラムが扱うアドレスはすべて仮想である．プログラムカウンタも，データ中に
格納されたポインタも，ロードとストアのためにプロセッサが計算するアドレスも
そうである．プログラムが使うのは，仮想アドレス空間のうちのいくつかの領域だけ
である（\cref{fig:l13-mapping} の左と右）．
\begin{itemize}
  \item 低位アドレスにある\emph{コード}とグローバルな\emph{データ}．
  \item その上にある\emph{ヒープ}．プログラムが動的にメモリを確保するにつれて
    高位アドレスの方向へ伸びる．
  \item アドレス空間の最上部にある\emph{スタック}．低位アドレスの方向へ伸びる．
\end{itemize}
ヒープとスタックの間には，プログラムが決して使わない膨大なアドレスの範囲が
ある．走っているプログラムはどれもこのように配置された自分専用の仮想アドレス
空間をもつので，異なるプログラムがまったく異なる内容に対して同じ仮想アドレスを
使うことになる．

\section{ページとページテーブル}
\label{sec:l13-pages}

仮想アドレスごとに対応づけを記録すると，アドレスの数だけテーブルのエントリが
必要になる．\source{Lesson 13，動画 6--7；H\&P §B.4．}そこで仮想記憶は，
キャッシュが個々のバイトではなくブロックを移動させるのと同じように，固定
サイズのアドレスの範囲を対応づける．

\begin{definition}[ページとページフレーム]
  仮想アドレス空間は等しい大きさの連続した範囲に分割され，その1つ1つを
  \term[ぺーじ]{ページ}[page]という．物理メモリも同じ大きさの範囲に分割され，
  その1つ1つを\term[ぺーじふれーむ]{ページフレーム}[page frame]，あるいは
  単にフレームという．その大きさは2のべき乗で，典型的には 4~KB であり，
  どのページもどのフレームにもちょうど収まる．
\end{definition}

\begin{definition}[ページテーブル]
  プログラムの\term[ぺーじてーぶる]{ページテーブル}[page table]は，その仮想
  アドレス空間のページごとに，そのページが物理メモリにあるかどうかと，ある
  場合はどのフレームがそれを保持しているかを記録する．
\end{definition}

オペレーティングシステムは，走っているプログラムごとに1つのページテーブルを
管理し，そのプログラムのページをフレームに置いたり取り除いたりするたびに
更新する．走っているプログラムのページテーブルは，オペレーティングシステムが
設定したレジスタを通じて見つけられ，そのプログラムのアドレスを変換するために
読まれる．読むのは通常プロセッサ自身であり，システムによっては
オペレーティングシステムである（\cref{sec:l13-tlb}）．\margin{このレジスタは
x86-64 では CR3，AArch64 では TTBR0/TTBR1 である．コンテキストスイッチが
これを入れ替えると，対応づけ全体が一度に変わる．}
\cref{fig:l13-mapping} は，2つのプログラムについて，こうして得られる
対応づけを示す．
\begin{itemize}
  \item どのページもどのフレームにでも置ける．仮想アドレス空間で隣り合う
    ページが物理メモリで隣り合っている必要はなく，異なるプログラムのフレームは
    互いに入り混じる．
  \item プログラムが決して使わないページは対応づけられず，フレームを占めない．
    物理メモリを消費するのは，アドレス空間のうち使われる領域だけである．
  \item 2つのプログラムのページテーブルが，同じフレームにページを対応づける
    こともある．これによりプログラムはメモリを共有でき，例えば同じプログラム
    が2つ走っているとき，そのコードを1つのコピーで済ませられる．
  \item 対応づけられたページがすべて物理メモリにある必要はない．残りは
    ディスクにある（\cref{sec:l13-page-faults}）．
\end{itemize}

\begin{figure}[htbp]
  \centering
  % The virtual address spaces of two programs mapped onto the frames of
% physical memory; one page of each program is on disk, and the code page
% is shared.  Page 0 (low addresses) is drawn at the bottom; page numbers
% are shown at the left inside each page.
\begin{tikzpicture}[x=1cm,y=1cm, font=\footnotesize\sffamily, >={Stealth[length=1.8mm]},
  pg/.style={draw, minimum width=1.5cm, minimum height=0.42cm, inner sep=0pt, font=\scriptsize\sffamily},
  num/.style={font=\scriptsize\sffamily, text=ln-muted, inner sep=1pt},
  ttl/.style={font=\scriptsize\sffamily\bfseries, text width=2.6cm, align=center, anchor=south, inner sep=1pt}]
  \def\h{0.42}
  \colorlet{lnpone}{ln-accent!15}
  \colorlet{lnptwo}{ln-tip!22}
  % ---------------- program 1 (left) and program 2 (right)
  \foreach \cx/\pcol in {0.75/lnpone, 9.05/lnptwo} {
    \node[pg, fill=\pcol] at (\cx,0.5*\h) {\iflangja{コード}{code}};
    \node[pg, fill=\pcol] at (\cx,1.5*\h) {\iflangja{データ}{data}};
    \node[pg, fill=\pcol] at (\cx,2.5*\h) {\iflangja{ヒープ}{heap}};
    \node[pg, fill=ln-box, minimum height=1.68cm, text=ln-muted] at (\cx,5*\h) {\iflangja{未使用}{unused}};
    \node[pg, fill=\pcol] at (\cx,7.5*\h) {\iflangja{スタック}{stack}};
    \foreach \i/\y in {0/0.5, 1/1.5, 2/2.5, {3--6}/6.5, 7/7.5} {
      \node[num, anchor=west] at (\cx-0.73,\y*\h) {\i};
    }
  }
  \node[ttl] at (0.75,8*\h+0.08) {\iflangja{プログラム 1 の\\仮想アドレス空間}{program 1:\\virtual address space}};
  \node[ttl] at (9.05,8*\h+0.08) {\iflangja{プログラム 2 の\\仮想アドレス空間}{program 2:\\virtual address space}};
  % ---------------- physical memory: frames 0..5 drawn from 2h to 8h
  \node[ttl] at (4.9,8*\h+0.08) {\iflangja{物理メモリ}{physical memory}};
  \node[pg, fill=lnpone] (f0) at (4.9,2.5*\h) {\iflangja{P1 データ}{P1 data}};
  \fill[lnpone] (4.15,3*\h) rectangle (4.9,4*\h);
  \fill[lnptwo] (4.9,3*\h) rectangle (5.65,4*\h);
  \node[pg] (f1) at (4.9,3.5*\h) {\iflangja{コード}{code}};
  \node[pg, fill=lnptwo] (f2) at (4.9,4.5*\h) {\iflangja{P2 データ}{P2 data}};
  \node[pg, text=ln-muted] (f3) at (4.9,5.5*\h) {\iflangja{空き}{free}};
  \node[pg, fill=lnpone] (f4) at (4.9,6.5*\h) {\iflangja{P1 スタック}{P1 stack}};
  \node[pg, fill=lnptwo] (f5) at (4.9,7.5*\h) {\iflangja{P2 ヒープ}{P2 heap}};
  % ---------------- disk below physical memory
  \draw[rounded corners=2pt, fill=ln-box] (3.45,-0.72) rectangle (6.35,0.12);
  \node[pg, fill=lnpone, minimum width=1.25cm] (d1) at (4.2,-0.3) {\iflangja{P1 ヒープ}{P1 heap}};
  \node[pg, fill=lnptwo, minimum width=1.25cm] (d2) at (5.6,-0.3) {\iflangja{P2 スタック}{P2 stack}};
  \node[num, anchor=north] at (4.9,-0.74) {\iflangja{ディスク}{disk}};
  % ---------------- mappings of program 1
  \draw[->] (1.5,0.5*\h) -- (f1.west);
  \draw[->] (1.5,1.5*\h) -- (f0.west);
  \draw[->] (1.5,7.5*\h) -- (f4.west);
  \draw[->, dashed] (0,2.5*\h) -- (-0.4,2.5*\h) -- (-0.4,-0.3) -- (d1.west);
  % ---------------- mappings of program 2
  \draw[->] (8.3,0.5*\h) -- (f1.east);
  \draw[->] (8.3,1.5*\h) -- (f2.east);
  \draw[->] (8.3,2.5*\h) -- (f5.east);
  \draw[->, dashed] (9.8,7.5*\h) -- (10.2,7.5*\h) -- (10.2,-0.3) -- (d2.east);
\end{tikzpicture}

  \caption{2つのプログラムの仮想アドレス空間（ページ 0--7，下が低位
    アドレス）を物理メモリの 6 つのフレームに対応づけたもの．両プログラムの
    ページ~0 は同じフレームに対応づけられている．各プログラムは 1 ページを
    ディスクにもち（破線），使われないページ 3--6 は対応づけられていない．}
  \label{fig:l13-mapping}
\end{figure}

ページテーブルは，使われるかどうかによらず仮想アドレス空間の\emph{すべての}
ページを記述する．最も単純な形（\cref{sec:l13-flat}）では，そのそれぞれに
エントリをもつ．したがって，記述するページの数を決めるのは仮想アドレス空間の
大きさとページの大きさであり，物理メモリの大きさは，1つのエントリが指せる
フレームの数を決めるにすぎない．\caution{ページテーブルの添字は仮想ページ
であって，フレームではない．物理メモリが少なければフレームは少なくなるが，
ページテーブルが記述すべきページの数は変わらない．}

\section{ページフォールト}
\label{sec:l13-page-faults}

走っているプログラム全体では，通常，物理メモリのフレーム数より多くのページが
対応づけられる．\source{Lesson 13，動画 8；H\&P §B.4．}収まりきらないページは
ディスクに置かれ，物理メモリは最近使われたページを保持する．

\begin{definition}[ページフォールト]
  ページテーブルがフレームに対応づけていないページへのアクセスは
  \term[ぺーじふぉーると]{ページフォールト}[page fault]，すなわち制御を
  オペレーティングシステムに移す例外を引き起こす．
\end{definition}

そのページがディスクにある場合，オペレーティングシステムは3つの手順で
ページフォールトを処理する．
\begin{enumerate}
  \item そのページのためにフレームを1つ選ぶ．空きフレームがなければ，最近
    使われていないページを追い出す．そのページが読み込まれてから書き込まれて
    いれば，先にディスクに書き戻す．追い出されたページのページテーブルは，
    そのページがもうメモリにないことを示すように更新する．
  \item フォールトしたページをディスクからそのフレームに読み込み，ページを
    このフレームに対応づけるようにページテーブルを更新する．
  \item フォールトを起こした命令を再実行する．今度はそのページがメモリに
    見つかる．これには正確な例外（第8章）が必要である．
\end{enumerate}

ディスクアクセスにはミリ秒，すなわち数百万クロックサイクルがかかるので，
ページフォールトはキャッシュミスよりはるかに高くつく．\margin{\SI{3}{\giga\hertz}
では，\SI{10}{\milli\second} のディスクシークは 3000 万サイクルであり，
\SI{100}{\micro\second} の NVMe SSD 読み出しでさえ 30 万サイクルである．}ページフォールトは
稀でなければならず，追い出すページを，オペレーティングシステムがソフトウェア
で入念なアルゴリズムを用いて選ぶだけの価値がある．

すべてのページフォールトがディスクを伴うわけではない．アドレスが，その
プログラムに割り当てられたどの領域にも属さない場合，読み込むものは何もなく，
オペレーティングシステムはそのアクセスを誤りとして扱い，通常はプログラムを
終了させる．\margin{ディスクに触れない場合はもう1つある．Linux は，ゼロクリアした
フレームを割り当てるだけ，あるいは共有ページを複製するだけで済むフォールトを
\emph{マイナー}フォールトと数え，記憶装置を読むものを\emph{メジャー}
フォールトとする．}

\begin{note}
  物理メモリは，ディスク上のページのキャッシュとして働く．ページは第12章の
  ブロックに，フレームはラインに，ページフォールトはキャッシュミスに，
  変更されたページだけをディスクに書き戻すことはダーティビットをもつ
  ライトバックキャッシュに対応する．どのページもどのフレームにでも置けるので
  配置はフルアソシアティブであるが，タグを探索する必要はない．ページテーブル
  はページを添字として直接引けるからである．
\end{note}

\section{アドレス変換}
\label{sec:l13-translation}

ページはページサイズの倍数から始まるので，仮想アドレスは，キャッシュ
アドレスのブロック番号とブロックオフセットとちょうど同じように，2つの
フィールドに分かれる．\source{Lesson 13，動画 9--10；H\&P §B.4．}

\begin{definition}[仮想ページ番号とページオフセット]
  $2^{k}$ バイトのページでは，仮想アドレスの下位 $k$ ビットが
  \term[ぺーじおふせっと]{ページオフセット}[page offset]，すなわちその
  バイトのページ内での位置を表す．残りの上位ビットが
  \term[かそうぺーじばんごう]{仮想ページ番号}[virtual page number]%
  \index{VPN|see{仮想ページ番号}}（VPN）であり，ページを識別する．
\end{definition}

フレームも同じように番号づけられる．物理アドレス $j \times 2^{k}$ から
始まるフレームの\term[ぶつりふれーむばんごう]{物理フレーム番号}[physical frame number]%
（PFN）\index{PFN|see{物理フレーム番号}}は~$j$ である．ページテーブルは VPN を
PFN に対応づけ，変換はオフセットを保ったまま前者を後者で置き換える
（\cref{fig:l13-translation}）．
\begin{equation}
  \begin{aligned}
    \text{VPN} &= \bigl\lfloor \mathit{VA} / 2^{k} \bigr\rfloor , &
    \text{オフセット} &= \mathit{VA} \bmod 2^{k} , \\
    \text{PFN} &= \text{ページテーブル}[\text{VPN}] , &
    \mathit{PA} &= \text{PFN} \times 2^{k} + \text{オフセット} ,
  \end{aligned}
  \label{eq:l13-translation}
\end{equation}
ここで $\mathit{VA}$ と $\mathit{PA}$ はそれぞれ仮想アドレスと物理アドレスで
ある．ページとそのフレームは同じ大きさであり，バイトはその中での位置を保つ
ので，オフセットはそのまま通り抜ける．2進表現では，物理アドレスは PFN の
ビットの後にオフセットのビットが続いたものである．\caution{変換はバイトを
そのページの外へ動かさないが，アラインされていないアクセスはページ境界を
またぐことがある．その2つの部分は別々のページに属し，2回の変換を要し，
離れたフレームに落ちることもある．}VPN は仮想アドレスより
$k$ ビット少なく，PFN は物理アドレスより $k$ ビット少ないので，仮想アドレスと
物理アドレスの幅が異なれば，この2つの番号の幅も異なる．そのページに対する
PFN がページテーブルにない場合，そのアクセスは代わりにページフォールトを
引き起こす．

\begin{figure}[htbp]
  \centering
  % Translation with a flat page table: 32-bit virtual address, 4 KB pages,
% 30-bit physical address.  The VPN indexes the page table, the PFN found
% there replaces it, and the page offset is copied unchanged.
\begin{tikzpicture}[x=1cm,y=1cm, font=\footnotesize\sffamily, >={Stealth[length=1.8mm]},
  bits/.style={font=\scriptsize\sffamily, text=ln-muted, inner sep=1pt},
  cellt/.style={font=\scriptsize\ttfamily}]
  % --- virtual address
  \draw[fill=ln-accent!15] (0,0) rectangle (3.2,0.55);
  \node at (1.6,0.275) {VPN\enspace\texttt{0x0040A}};
  \draw[fill=white] (3.2,0) rectangle (5.6,0.55);
  \node at (4.4,0.275) {\iflangja{オフセット}{offset}\enspace\texttt{0x7C6}};
  \node[bits, anchor=south] at (1.6,0.55) {31\,\dots\,12};
  \node[bits, anchor=south] at (4.4,0.55) {11\,\dots\,0};
  \node[anchor=west] at (5.7,0.275) {\iflangja{仮想アドレス}{virtual address}};
  % --- page table: V | PFN, 6 drawn rows
  \node[font=\sffamily\scriptsize\bfseries, anchor=south] at (2.1,-0.92) {\iflangja{ページテーブル}{page table}};
  \node[bits, anchor=south] at (1.15,-1.2) {V};
  \node[bits, anchor=south] at (2.35,-1.2) {PFN};
  \fill[ln-accent!15] (0.9,-2.46) rectangle (3.3,-2.88);
  \foreach \r/\v/\pfn/\lab in {0/1/0x00A51/0, 1/0/{{\sffamily ---}}/1, 2/{}/{}/{}, 3/1/0x1F3C5/0x0040A, 4/{}/{}/{}, 5/1/0x3B207/0xFFFFF} {
    \draw (0.9,-1.2-0.42*\r) rectangle (1.4,-1.62-0.42*\r);
    \draw (1.4,-1.2-0.42*\r) rectangle (3.3,-1.62-0.42*\r);
    \node[cellt] at (1.15,-1.41-0.42*\r) {\v};
    \node[cellt] at (2.35,-1.41-0.42*\r) {\pfn};
    \node[bits, font=\scriptsize\ttfamily, anchor=west] at (3.35,-1.41-0.42*\r) {\lab};
  }
  \foreach \r in {2,4} {
    \node[font=\small, text=ln-muted, inner sep=0pt] at (1.15,-1.41-0.42*\r) {$\vdots$};
    \node[font=\small, text=ln-muted, inner sep=0pt] at (2.35,-1.41-0.42*\r) {$\vdots$};
  }
  % --- VPN selects an entry
  \draw[->] (0.3,0) -- (0.3,-2.67) -- (0.9,-2.67);
  % --- physical address
  \draw[->] (2.35,-3.72) -- (2.35,-4.3);
  \draw[fill=ln-accent!15] (0.8,-4.85) rectangle (3.2,-4.3);
  \node at (2.0,-4.575) {PFN\enspace\texttt{0x1F3C5}};
  \draw[fill=white] (3.2,-4.85) rectangle (5.6,-4.3);
  \node at (4.4,-4.575) {\iflangja{オフセット}{offset}\enspace\texttt{0x7C6}};
  \node[bits, anchor=north] at (2.0,-4.85) {29\,\dots\,12};
  \node[bits, anchor=north] at (4.4,-4.85) {11\,\dots\,0};
  \node[anchor=west] at (5.7,-4.575) {\iflangja{物理アドレス}{physical address}};
  % --- offset is copied
  \draw[->] (5.0,0) -- (5.0,-4.3);
\end{tikzpicture}

  \caption{4~KB のページと 30 ビットの物理アドレスのもとでの仮想アドレス
    \texttt{0x0040A7C6} の変換（\cref{ex:l13-translation}）．VPN が
    ページテーブルエントリを選び，その PFN が VPN を置き換え，オフセットは
    そのまま写される．V~$=0$ のエントリであればページフォールトとなる．}
  \label{fig:l13-translation}
\end{figure}

\begin{example}
  \label{ex:l13-translation}
  4~KB のページ（$k = 12$）をもつ 32 ビットプロセッサが 1~GB の物理メモリを
  もつとすると，物理アドレスは 30 ビットである．ページオフセットはビット
  11--0，VPN はビット 31--12（20 ビット）であり，PFN は 18 ビットである．
  仮想アドレス \texttt{0x0040A7C6} の VPN は \texttt{0x0040A}，オフセットは
  \texttt{0x7C6} である．\cref{fig:l13-translation} のページテーブルはページ
  \texttt{0x0040A} をフレーム \texttt{0x1F3C5} に対応づけるので，物理アドレス
  は $\texttt{0x1F3C5} \times 2^{12} + \texttt{0x7C6} = \texttt{0x1F3C57C6}$
  である．仮想アドレス \texttt{0x000017FF} はページ \texttt{0x00001} に属し，
  同じページテーブルはこれをメモリにないと記している．このアドレスへの
  アクセスはページフォールトを引き起こす．
\end{example}

\tip{4~KB のページで 16 進表記のアドレスを使うと，下 3 桁がオフセットである．
変換はその手前の桁を置き換える．}

\begin{keypoint}
  変換は仮想ページ番号を物理フレーム番号に置き換え，ページオフセットを
  そのまま残す．フレーム番号を与えるのは，オペレーティングシステムが管理する
  ページテーブルである．そのページに対するフレーム番号がなければ，そのアクセス
  はページフォールトを引き起こす．
\end{keypoint}

\section{単一レベルページテーブルの大きさ}
\label{sec:l13-flat}

最も単純なページテーブルは，1つの配列である．\source{Lesson 13，動画
11--12；H\&P §B.4．}

\begin{definition}[単一レベルページテーブル]
  \term[たんいつれべるぺーじてーぶる]{単一レベルページテーブル}[flat page table]%
  とは，仮想ページごとに1つのエントリをもち，VPN で直接添字づけされる配列で
  ある．あるページに対応するその要素を，そのページの
  \term[ぺーじてーぶるえんとり]{ページテーブルエントリ}[page table entry]%
  \index{PTE|see{ページテーブルエントリ}}（PTE）という．
\end{definition}

PTE は，そのページの PFN といくつかの状態ビット，すなわちページが物理メモリに
あるときだけセットされる有効ビットと，そのページにどのようにアクセスしてよいか
（読み出し，書き込み，実行）を記録するビットを保持する．\margin{x86-64 の PTE は
8 バイトで，PFN のほかに存在ビット，読み書き，ユーザ/スーパバイザ，アクセス
ビット，ダーティビット，そして最上位の実行禁止ビットをもつ．}PTE の大きさは典型的に
はアドレスと同程度，32 ビットプロセッサで 4 バイト，64 ビットプロセッサで 8
バイトであり，PFN と状態ビットを収めるだけの余地がある．どの仮想ページにも
エントリがあるので，
\begin{equation}
  \begin{aligned}
    &\text{単一レベルページテーブルの大きさ} \\
    &\qquad =\; \frac{\text{仮想アドレス空間の大きさ}}{\text{ページサイズ}}
      \times \text{PTE の大きさ} .
  \end{aligned}
  \label{eq:l13-flat-size}
\end{equation}
\cref{eq:l13-flat-size} に現れない量が2つある．物理メモリの大きさは，
エントリのうち PFN が占めるビット数に影響するだけである．プログラムが実際に
使うメモリの量はまったく関係しない．使われないページにもエントリがあるから
である．走っているプログラムはどれも，この大きさのページテーブルを必要と
する．\tip{指数で考えるとよい．$n$ ビットのアドレス空間，$2^{k}$ バイトの
ページ，$2^{s}$ バイトのエントリなら，テーブルは $2^{\,n-k+s}$ バイトである．
ページのビットを引き，エントリのビットを足す．}

\begin{example}
  \label{ex:l13-flat-size}
  32 ビットの仮想アドレス，4~KB のページ，4 バイトのエントリでは，単一レベル
  ページテーブルは $2^{32}/2^{12} = 2^{20}$ 個のエントリをもち，走っている
  プログラム1つにつき $2^{20} \times 4$ バイト $= 4$~MB を占める．数キロバイト
  しか使わないプログラムでも同じである．64 ビットの仮想アドレス，4~KB の
  ページ，8 バイトのエントリなら $2^{52}$ 個のエントリをもち，$2^{55}$ バイト
  （32~PB），すなわちどのメモリにもディスクにも収まらない量を占めることに
  なる．64 ビットのアドレスのうち 48 ビットしか使わない場合でも，単一レベル
  ページテーブルは $2^{36} \times 8 = 2^{39}$ バイト，すなわちプログラム
  1つにつき 512~GB を占める．
\end{example}

\section{多段ページテーブル}
\label{sec:l13-multilevel}

単一レベルページテーブルが大きいのは，プログラムがアドレス空間の下端と上端
しか使わないにもかかわらず，空間全体に対するエントリをもつからである．
\source{Lesson 13，動画 13--17；H\&P §B.4--B.5．}多段ページテーブルは，
その間の使われない範囲に対するエントリを保持しないようにする．

\begin{definition}[多段ページテーブル]
  \term[ただんぺーじてーぶる]{多段ページテーブル}[multi-level page table]とは，
  小さなテーブルからなる木である．VPN は段ごとに1つのフィールドに分割される．
  最初のフィールドは外側ページテーブルに添字を与え，そのエントリは次の段の
  テーブルを指す．最後の段のテーブル，すなわち内側ページテーブルのエントリが
  PFN を保持する．担当するアドレス空間の部分がまったく使われていないエントリは
  無効と記され，その下にはテーブルが存在しない．
\end{definition}

2段の場合，仮想アドレスは外側の添字 $p_1$，内側の添字 $p_2$，ページ
オフセットに分かれ（\cref{fig:l13-two-level}），変換は2つの手順を踏む．
\begin{enumerate}
  \item 外側ページテーブルのエントリ $p_1$ が，内側ページテーブルの位置を
    与える．
  \item その内側ページテーブルのエントリ $p_2$ が PFN を与え，これが
    \cref{eq:l13-translation} のようにオフセットと組み合わされる．
\end{enumerate}
どちらかのエントリが無効なら PFN は得られず，そのアクセスはページフォールトを
引き起こす．テーブルの大きさは，通常，1つがちょうど1ページに収まるように
定められる．

\begin{figure}[htbp]
  \centering
  % Two-level translation: 32-bit virtual address split into a 10-bit outer
% index p1, a 10-bit inner index p2 and a 12-bit offset.  Outer entry p1
% points to an inner table, whose entry p2 holds the PFN.
\begin{tikzpicture}[x=1cm,y=1cm, font=\footnotesize\sffamily, >={Stealth[length=1.8mm]},
  bits/.style={font=\scriptsize\sffamily, text=ln-muted, inner sep=1pt},
  cellt/.style={font=\scriptsize\ttfamily},
  ttl/.style={font=\sffamily\scriptsize\bfseries, anchor=south, inner sep=1pt}]
  % --- virtual address
  \draw[fill=ln-accent!15] (0,0) rectangle (2.4,0.55);
  \node at (1.2,0.275) {$p_1 = \texttt{1}$};
  \draw[fill=ln-box] (2.4,0) rectangle (4.8,0.55);
  \node at (3.6,0.275) {$p_2 = \texttt{0x00A}$};
  \draw[fill=white] (4.8,0) rectangle (7.2,0.55);
  \node at (6.0,0.275) {\iflangja{オフセット}{offset}\enspace\texttt{0x7C6}};
  \node[bits, anchor=south] at (1.2,0.55) {31\,\dots\,22};
  \node[bits, anchor=south] at (3.6,0.55) {21\,\dots\,12};
  \node[bits, anchor=south] at (6.0,0.55) {11\,\dots\,0};
  \node[anchor=west] at (7.3,0.275) {\iflangja{仮想アドレス}{virtual address}};
  % --- outer table: V | inner table
  \node[ttl] at (1.7,-0.92) {\iflangja{外側テーブル}{outer table}};
  \node[bits, anchor=south] at (1.1,-1.2) {V};
  \node[bits, anchor=south] at (1.9,-1.2) {\iflangja{内側}{inner}};
  \fill[ln-accent!15] (0.9,-1.62) rectangle (2.5,-2.04);
  \foreach \r/\v/\p/\lab in {0/1/A/0, 1/1/B/{}, 2/1/C/2, 3/0/{{\sffamily ---}}/3, 4/{}/{}/{}, 5/1/D/1023} {
    \draw (0.9,-1.2-0.42*\r) rectangle (1.3,-1.62-0.42*\r);
    \draw (1.3,-1.2-0.42*\r) rectangle (2.5,-1.62-0.42*\r);
    \node[cellt] at (1.1,-1.41-0.42*\r) {\v};
    \node[cellt] at (1.8,-1.41-0.42*\r) {\p};
    \node[bits, anchor=east] at (0.85,-1.41-0.42*\r) {\lab};
  }
  \node[font=\small, text=ln-muted, inner sep=0pt] at (1.1,-1.41-0.42*4) {$\vdots$};
  \node[font=\small, text=ln-muted, inner sep=0pt] at (1.9,-1.41-0.42*4) {$\vdots$};
  \draw[->] (0.15,0) -- (0.15,-1.83) -- (0.9,-1.83);
  % index of outer entry 1, above the p1 arrow that enters its row
  \node[bits, anchor=south east] at (0.85,-1.78) {1};
  % --- inner table B: V | PFN, top edge at the height of outer entry 1
  \node[ttl] at (4.35,-1.55) {\iflangja{内側テーブル B}{inner table B}};
  \node[bits, anchor=south] at (3.65,-1.83) {V};
  \node[bits, anchor=south] at (4.6,-1.83) {PFN};
  \fill[ln-accent!15] (3.4,-3.09) rectangle (5.3,-3.51);
  \foreach \r/\v/\pfn/\lab in {0/1/0x2A017/0, 1/0/{{\sffamily ---}}/1, 2/{}/{}/{}, 3/1/0x1F3C5/10, 4/{}/{}/{}, 5/1/0x0C4E2/1023} {
    \draw (3.4,-1.83-0.42*\r) rectangle (3.9,-2.25-0.42*\r);
    \draw (3.9,-1.83-0.42*\r) rectangle (5.3,-2.25-0.42*\r);
    \node[cellt] at (3.65,-2.04-0.42*\r) {\v};
    \node[cellt] at (4.6,-2.04-0.42*\r) {\pfn};
    \node[bits, anchor=east] at (3.35,-2.04-0.42*\r) {\lab};
  }
  \foreach \r in {2,4} {
    \node[font=\small, text=ln-muted, inner sep=0pt] at (3.65,-2.04-0.42*\r) {$\vdots$};
    \node[font=\small, text=ln-muted, inner sep=0pt] at (4.6,-2.04-0.42*\r) {$\vdots$};
  }
  % --- pointer from outer entry 1 to the start of inner table B
  \fill (2.3,-1.83) circle (1pt);
  \draw[->] (2.3,-1.83) -- (3.4,-1.83);
  % --- p2 selects an entry of B
  \draw[->] (4.4,0) -- (4.4,-0.35) -- (5.75,-0.35) -- (5.75,-3.30) -- (5.3,-3.30);
  % --- physical address
  \draw[->] (4.6,-4.35) -- (4.6,-5.0);
  \draw[fill=ln-accent!15] (2.4,-5.55) rectangle (4.8,-5.0);
  \node at (3.6,-5.275) {PFN\enspace\texttt{0x1F3C5}};
  \draw[fill=white] (4.8,-5.55) rectangle (7.2,-5.0);
  \node at (6.0,-5.275) {\iflangja{オフセット}{offset}\enspace\texttt{0x7C6}};
  \node[bits, anchor=north] at (3.6,-5.55) {29\,\dots\,12};
  \node[bits, anchor=north] at (6.0,-5.55) {11\,\dots\,0};
  \node[anchor=west] at (7.3,-5.275) {\iflangja{物理アドレス}{physical address}};
  % --- offset is copied
  \draw[->] (6.8,0) -- (6.8,-5.0);
\end{tikzpicture}

  \caption{仮想アドレス \texttt{0x0040A7C6} の2段の変換
    （\cref{ex:l13-two-level}）．外側のエントリ~1 は内側テーブル~B を指す．
    A，C，D のエントリは図示していない内側テーブルを指す．エントリ~3 のような
    無効な外側エントリには，内側テーブルがそもそも存在しない．}
  \label{fig:l13-two-level}
\end{figure}

\begin{example}
  \label{ex:l13-two-level}
  \cref{ex:l13-translation} の 20 ビットの VPN を，10 ビットの外側の添字
  （ビット 31--22）と 10 ビットの内側の添字（ビット 21--12）に分割する．
  4 バイトのエントリなら，どのテーブルも $2^{10}$ 個のエントリをもち 4~KB，
  すなわちちょうど1ページを占める．仮想アドレス \texttt{0x0040A7C6} の VPN は
  $\texttt{0x0040A} = 0000000001\,0000001010_2$ なので，$p_1 = 1$，
  $p_2 = \texttt{0x00A} = 10$ である．外側のエントリ~1 は内側テーブル~B を
  指し，B のエントリ~10 が PFN \texttt{0x1F3C5} を保持する．物理アドレスは
  単一レベルページテーブルの場合と同じ \texttt{0x1F3C57C6} である．多段
  ページテーブルは，同じ対応づけを別の形で格納しているにすぎない．
\end{example}

\subsection{多段ページテーブルの大きさ}

外側ページテーブルは常に存在するが，内側ページテーブルが存在するのは，
その担当するアドレス空間の部分で少なくとも1つのページが使われている場合だけ
である．\cref{ex:l13-two-level} の構成では，内側テーブル1つが $2^{10}$
ページ，すなわち仮想アドレス空間の 4~MB を担当する．

\begin{example}
  \label{ex:l13-two-level-size}
  あるプログラムが，32 ビットのアドレス空間の下端の 10~MB をコード，データ，
  ヒープに（アドレス \texttt{0x00000000}--\texttt{0x009FFFFF}），上端の 1~MB を
  スタックに（\texttt{0xFFF00000}--\texttt{0xFFFFFFFF}）使うとする．
  \cref{ex:l13-two-level} の2段ページテーブルでは，下端の 10~MB はそれぞれ
  4~MB を担当する外側エントリ 0，1，2 の範囲に，スタックは外側エントリ 1023 の
  範囲に入る．ページテーブルは外側テーブル1つと内側テーブル4つからなり，
  単一レベルページテーブルの 4~MB ではなく
  $5 \times 4~\text{KB} = 20$~KB である．\sidenote{同じ問題への別の答えが\emph{逆ページテーブル}である．ページごと
  ではなくフレームごとに1エントリをもち，VPN のハッシュで引く．大きさは物理
  メモリで決まりアドレス空間とともに増えることはないが，衝突すると連鎖を
  たどる必要があり，ページの共有にも別の仕掛けがいる．PowerPC や UltraSPARC
  が用いた（H\&P §B.4）．}
\end{example}

2段ページテーブルが単一レベルのものより大きくなるのは，すべての外側エントリに
内側テーブルが必要な場合，すなわち 4~MB の 1024 個の範囲のそれぞれで少なくとも
1つのページが使われている場合だけである．そのとき
$4~\text{KB} + 1024 \times 4~\text{KB} \approx 4.004$~MB を占め，単一レベルの
テーブルより外側テーブルの分だけ大きくなる．実際のプログラムはアドレス空間の
連続したいくつかの領域を使うので，必要な内側テーブルはわずかである．

\subsection{64 ビットアドレスのためのより多くの段}

64 ビットのアドレス空間には2段では足りない．8 バイトのエントリからなる 4~KB
のテーブルは $2^{9}$ 個のエントリをもつので，1ページに収まるテーブルからなる
各段が消費できるのは VPN の 9 ビットだけである．48 ビットの仮想アドレスでは
VPN は 36 ビットである．9 ビットの内側テーブルをもつ2段ページテーブルなら
27 ビットの外側の添字が必要で，外側テーブルはプログラム1つにつき
$2^{27} \times 8$ バイト $= 1$~GB になる．代わりに 36 ビットを 9 ビットずつ
4つのフィールドに分割すると，どのテーブルも1ページに収まる4段ページテーブルが
得られる．x86-64 プロセッサはこの構成を用いる．段を1つ増やすごとに，仮想
アドレスをさらに 9 ビット分収められるようになる．\margin{x86-64 は4つの段を
PML4，PDPT，PD，PT と呼ぶ．Ice Lake 以降のプロセッサは，57 ビットすなわち
128~PB のアドレスのために5段目（\texttt{CR4.LA57}）を備える．}

\begin{example}
  \label{ex:l13-four-level}
  このような4段ページテーブルでは，最後の段のテーブル1つが $2^{9}$ ページ，
  すなわちアドレス空間の 2~MB を担当し，第3段のテーブルはその $2^{9}$ 倍の
  1~GB を，第2段のテーブルは 512~GB を，そして第1段のただ1つのテーブルは
  $2^{48}$ バイト（256~TB）すべてを担当する．アドレス空間の下端で 8~MB，
  上端で 8~MB を使うプログラムは，それぞれの端について最後の段のテーブル4つと
  第3段および第2段のテーブルを1つずつ必要とし，それらはすべて第1段のただ1つの
  テーブルの下にぶら下がる．合計で
  $1 + 2 \times (4 + 1 + 1) = 13$ 個のテーブル，すなわち
  $13 \times 4~\text{KB} = 52$~KB である．単一レベルページテーブルなら
  512~GB を占めることになる（\cref{ex:l13-flat-size}）．
\end{example}

\begin{keypoint}
  ページテーブルは仮想アドレス空間全体を記述しなければならないが，多段ページ
  テーブルが記憶領域を必要とするのは，使われている部分についてだけである．
  その代償は，変換のたびに段の数だけテーブルを読むことである．
\end{keypoint}

\section{ページサイズの選択}
\label{sec:l13-page-size}

ページをどれだけ大きくすべきかは，2つの効果の折り合いで決まる．
\source{Lesson 13，動画 18；H\&P §B.4．}
\begin{itemize}
  \item \textbf{ページが大きい}ほどページの数は少なくなる．PTE が減って
    ページテーブルが小さくなり，さらに \cref{sec:l13-tlb} で示すように，
    プロセッサが保持する変換1つが覆うメモリが増える．
  \item \textbf{ページが小さい}ほど，部分的にしか使われないページで無駄になる
    メモリが減る．
\end{itemize}
1つのページと1つのフレームのバイト数である
\term[ぺーじさいず]{ページサイズ}[page size]は，この両者の均衡をとらなければ
ならない．

\begin{definition}[内部断片化]
  メモリはページ単位で対応づけられる．したがって，大きさがページサイズの倍数
  でない領域は，最後のページの一部を使わないまま残し，その部分は他のどの
  プログラムも使えない．このようにして失われるメモリを
  \term[ないぶだんぺんか]{内部断片化}[internal fragmentation]という．
\end{definition}

大きさが任意の領域は平均して半ページを使わないまま残すので，内部断片化は
ページサイズに比例して増える．\caution{ページングは\emph{外部断片化}，すなわち大きさの異なる割り当ての間に
残る無駄を取り除く．どのページもどのフレームにも収まるからである．残るのは
最後のページの内側，内部断片化の方である．}

\begin{example}
  \label{ex:l13-page-size}
  4 バイトの PTE をもつ 32 ビットプロセッサ上のあるプログラムが，3つの領域
  （コードとデータ，ヒープ，スタック）をもち，そのそれぞれが部分的にしか
  使われないページで終わっているとする．\cref{tab:l13-page-size} は3つの
  ページサイズを比較したものである．4~KB のページから 4~MB のページにすると，
  単一レベルページテーブルは $2^{10}$ 分の1，すなわち 4~MB から 4~KB に
  縮むが，内部断片化の期待値は同じ倍率で，6~KB から 6~MB に増える．
\end{example}

\begin{table}[htbp]
  \centering\small
  \caption{\cref{ex:l13-page-size} のプログラムについての，ページテーブルの
    大きさと内部断片化の期待値（領域当たり半ページ）．}
  \label{tab:l13-page-size}
  \begin{tabular}{@{}rrrr@{}}
    \toprule
    ページサイズ & ページ数 & 単一レベルページテーブル & 内部断片化 \\
    \midrule
    4~KB  & $2^{20}$ & 4~MB   & $3 \times 2~\text{KB} = 6$~KB \\
    64~KB & $2^{16}$ & 256~KB & $3 \times 32~\text{KB} = 96$~KB \\
    4~MB  & $2^{10}$ & 4~KB   & $3 \times 2~\text{MB} = 6$~MB \\
    \bottomrule
  \end{tabular}
\end{table}

ページはキャッシュブロック（第12章）よりはるかに大きい．ページごとに PTE が
必要であり，ディスクから持ち込むページごとにディスクアクセスがかかるから
である．通常のページサイズは 4~KB である．\margin{x86 では
1985 年の 386 以来 4~KB がページサイズである．AArch64 では，オペレーティング
システムが 4，16，64~KB の粒度から選べる．}プロセッサは，x86-64 の 2~MB や
1~GB のように，はるかに大きなページも備えている．オペレーティングシステムは，
プログラムが完全に埋める大きな領域にこれを使うことができ，その場合，内部
断片化はほとんど問題にならない．

\section{変換を伴うメモリアクセス時間}
\label{sec:l13-access-time}

プログラムが行うメモリアクセスは，命令フェッチであれロードやストアであれ，
すべて仮想アドレスを使うので，物理メモリにアクセスする前に変換されなければ
ならない．\source{Lesson 13，動画 19--21；H\&P §B.4．}ページテーブル自体も
メモリに置かれているので，変換はメモリ読み出しからなる．2段ページテーブルの
場合，プロセッサは \texttt{LW R1,4(R2)} のようなロードを4つの手順で実行する．
\begin{enumerate}
  \item R2 に 4 を加えて仮想アドレスを求め，その VPN を取り出す．
  \item 外側の PTE の物理アドレスを（外側テーブルのアドレスに $p_1$ と PTE の
    大きさの積を加えて）計算し，その PTE を読む．
  \item 同じようにして，外側の PTE にあるテーブルのアドレスと $p_2$ から内側の
    PTE のアドレスを計算し，その PTE を読む．
  \item PFN にオフセットを付け加え，この物理アドレスにあるデータにキャッシュを
    通じてアクセスする．
\end{enumerate}
PTE の読み出しはいずれも通常のメモリアクセスであり，キャッシュを通り，ヒット
することもミスすることもある．PTE がキャッシュに保持されていなければ，1回
ごとに主記憶への完全なアクセスがかかる．\margin{x86 プロセッサは通常の
キャッシュとは別に，上位段のエントリを専用の\emph{ページング構造キャッシュ}
に保持する．そのためウォークはたいてい最下段の読み出しだけで済む．
Intel SDM 第3巻．}$L$ 段のページテーブルでは，プログラム
のアクセス1回が $L + 1$ 回のメモリアクセスになる．PTE の読み出しがデータ
アクセスと同じ平均メモリアクセス時間（AMAT）をもつとすれば，
\begin{equation}
  \text{変換を伴うアクセス時間} \;\approx\; (L + 1) \times \text{AMAT} .
  \label{eq:l13-walk-time}
\end{equation}

\begin{example}
  \label{ex:l13-access-time}
  あるキャッシュのヒット時間が 2 サイクル，ミス率が \SI{5}{\percent}，
  ミスペナルティが 40 サイクルなら，$\text{AMAT} = 2 + 0.05 \times 40 = 4$
  サイクルである（\cref{eq:l12-amat}）．変換がなければメモリアクセスは平均
  4 サイクルだが，単一レベルページテーブルでは $2 \times 4 = 8$ サイクル，
  4段ページテーブルでは $5 \times 4 = 20$ サイクルになる．個々のアクセスの
  ばらつきは大きい．上位3段の PTE がヒットし（$3 \times 2$ サイクル），最後の
  段の PTE がミスし（$2 + 40$ サイクル），データがヒットする（2 サイクル）
  ロードは 50 サイクルかかり，これは変換がなければキャッシュヒットで済んだ
  時間の 25 倍である．
\end{example}

上位の段のわずかなエントリはどの変換でも読まれるので，たいていキャッシュに
保持される．数の多い最後の段のエントリはより頻繁にミスする．いずれにせよ，
より速く変換する方法がなければ，仮想記憶はすべての命令フェッチ，ロード，
ストアを大きな倍率で遅くしてしまう．

\section{TLB}
\label{sec:l13-tlb}

あるページにアクセスしたプログラムは，近いうちに同じページに再びアクセスする
可能性が高いので，ある時点で必要な変換はわずかであり，それが何度も繰り返し
使われる．\source{Lesson 13，動画 22--24；H\&P §B.4．}そこでプロセッサは，
それをキャッシュする．

\begin{definition}[TLB]
  \term{TLB}[translation lookaside buffer]とは，最近の変換を保持する小容量で
  高速なキャッシュである．各エントリは，有効ビット，VPN から取ったタグ
  （フルアソシアティブの TLB では VPN 全体），そのページが対応づけられている
  PFN，およびその PTE のアクセスビットを保持する．アドレスの VPN が有効な
  エントリと一致すれば，PFN はそのエントリから得られ，ページテーブルは
  読まれない．
\end{definition}

TLB は第12章の意味でのキャッシュであるが，データキャッシュよりはるかに小さく
できる．\tip{TLB が
一度に覆えるメモリをその\emph{リーチ}という．エントリ数 $\times$ ページ
サイズであり，4~KB の 64 エントリなら 256~KB，同じ 64 エントリでも
2~MB のページなら 128~MB になる．}理由は3つある．
\begin{itemize}
  \item 格納するのは変換だけでデータではない．1エントリは数バイトで済む．
  \item 1つのエントリが，1つのキャッシュブロックではなく1ページ全体に対応する
    ので，わずかなエントリでプログラムが現に使っているメモリを覆える．
  \item 格納するのは変換の最終結果であって，個々の段のエントリではない．
    ページテーブルが何段であろうと，TLB ヒットは1回の探索で済む．
\end{itemize}

\begin{example}
  \label{ex:l13-tlb-size}
  64 バイトのブロックをもつ 128~KB のデータキャッシュには 2048 本のラインが
  ある．4~KB のページなら，$128/4 = 32$ 個の TLB エントリでキャッシュが
  保持するのと同じだけのメモリを覆える．キャッシュされたブロックがすべて
  別々のページにある場合には 2048 個が必要になるが，空間的局所性により実際に
  必要な数はこの上限をはるかに下回る．
\end{example}

\subsection{TLB ミス}

アドレスの VPN に一致する有効な TLB エントリがないとき，
\term[TLB みす]{TLB ミス}[TLB miss]が生じる．このとき変換はページテーブルから
得なければならない．

\begin{definition}[ページテーブルウォーク]
  \term[ぺーじてーぶるうぉーく]{ページテーブルウォーク}[page table walk]とは，
  \cref{sec:l13-multilevel} で述べたように，あるアドレスのページテーブル
  エントリを，外側ページテーブルから PFN を保持するエントリまで段ごとに読む
  ことである．
\end{definition}

TLB ミスの後，ウォークが PFN を見つけ，その変換が TLB に書き込まれ（TLB が
満杯なら既存のエントリを置き換え），アクセスが進む
（\cref{fig:l13-tlb-flow}）．\caution{TLB ミスはページフォールトではない．
ウォークはたいてい有効な PTE を見つける．}

TLB はページテーブルと一致していなければならない．オペレーティングシステムが
対応づけを変更したとき，例えばページを追い出したときには，そのページの TLB
エントリを無効化する．また，プロセッサが別のプログラムに切り替わるときには，
前のプログラムのエントリを無効化するか，プログラムの識別子をタグとして
付けておく．\sidenote{このタグを x86 では PCID，Arm では ASID という．これがあれば
コンテキストスイッチで TLB を捨てずに済む．難しいのは1つの対応づけを無効化
する場合である．マルチコアプロセッサでは，そのエントリを保持しうるすべての
コアがそれを捨てなければならないので，オペレーティングシステムはプロセッサ
間割り込みを送り，その応答を待つ．この \emph{TLB シュートダウン} のために，
メモリの対応づけを解除する費用は対応づける費用よりはるかに大きい．}

\begin{figure}[htbp]
  \centering
  % Translation of one memory access: TLB hit, TLB miss with a page table
% walk, and page fault handled by the operating system.
\begin{tikzpicture}[x=1cm,y=1cm, font=\footnotesize\sffamily, >={Stealth[length=1.8mm]},
  stp/.style={draw, fill=ln-box, align=center, inner sep=3pt, minimum height=0.8cm, text width=3.0cm},
  os/.style={stp, fill=white},
  dec/.style={draw, diamond, aspect=2.4, fill=ln-accent!15, align=center, inner sep=1pt},
  lab/.style={font=\sffamily\scriptsize, text=ln-muted}]
  % --- column 1: TLB lookup and the access itself
  \node[stp] (split) at (0,0)
    {\iflangja{仮想アドレスを VPN と\\オフセットに分ける}{split the virtual address\\into VPN and offset}};
  \node[stp] (look) at (0,-1.3)
    {\iflangja{VPN で TLB を検索}{look up the VPN\\in the TLB}};
  \node[dec] (hit) at (0,-2.65) {\iflangja{TLB ヒット?}{TLB hit?}};
  \node[stp] (acc) at (0,-5.45)
    {\iflangja{物理アドレス $=$ PFN と\\オフセット; データに\\アクセス}{physical address $=$\\PFN and offset;\\access the data}};
  % --- column 2: page table walk
  \node[stp] (walk) at (3.9,-2.65)
    {\iflangja{ページテーブルウォーク\\（ハードウェアまたは OS）}{page table walk\\(hardware or OS)}};
  \node[dec] (valid) at (3.9,-4.0) {\iflangja{PTE 有効?}{PTE valid?}};
  \node[stp] (fill) at (3.9,-5.45)
    {\iflangja{VPN と PFN を\\TLB に書き込む}{write VPN and PFN\\into the TLB}};
  % --- column 3: page fault
  \node[os] (fault) at (7.8,-4.0)
    {\iflangja{ページフォールト（OS）:\\ページをディスクから\\読み込み，PTE を更新}{page fault (OS):\\load the page from disk,\\update the page table}};
  % --- arrows
  \draw[->] (split) -- (look);
  \draw[->] (look) -- (hit);
  \draw[->] (hit) -- node[lab, right, pos=0.12] {\iflangja{はい}{yes}} (acc);
  \draw[->] (hit) -- node[lab, above] {\iflangja{いいえ}{no}} (walk);
  \draw[->] (walk) -- (valid);
  \draw[->] (valid) -- node[lab, right, pos=0.2] {\iflangja{はい}{yes}} (fill);
  \draw[->] (fill) -- (acc);
  \draw[->] (valid) -- node[lab, above] {\iflangja{いいえ}{no}} (fault);
  \draw[->] (fault.north) -- (7.8,-1.3) -- node[lab, above, pos=0.5] {\iflangja{命令を再実行}{restart the instruction}} (look.east);
\end{tikzpicture}

  \caption{メモリアクセスにおける変換．TLB ヒットならただちに PFN が得られ，
    TLB ミスならページテーブルウォークが必要になる．PTE が無効であれば
    ページフォールトとなり，ページがメモリに入った後に命令が再実行される．}
  \label{fig:l13-tlb-flow}
\end{figure}

ウォークは2通りの方法で行える．
\begin{description}
  \item[ソフトウェアによる方法] TLB ミスが例外を起こし，オペレーティング
    システム内のハンドラがページテーブルをたどって変換を TLB に書き込む．
    ハードウェアは単純で済み，ページテーブルをどのように構成してもよいが，
    TLB ミスのたびに例外の代償を払う．MIPS プロセッサは TLB ミスをこの方法で
    処理する．
  \item[ハードウェアによる方法] プロセッサがページテーブルウォーカ，すなわち
    PTE を読んで自ら TLB を埋める論理回路を内蔵する．TLB ミスの処理は
    はるかに速いが，ページテーブルはハードウェアが想定する形式でなければ
    ならない．x86 系のような高性能プロセッサは，ページテーブルウォークを
    ハードウェアで行う．
\end{description}

\needspace{6\baselineskip}
\begin{keypoint}
  どのメモリアクセスにも変換が必要である．TLB は，よくある場合である TLB
  ヒットを1回の高速な探索で済ませる．ページテーブルウォークは，ハードウェア
  によるものであれオペレーティングシステムによるものであれ TLB ミスのときだけ
  必要であり，ページフォールトハンドラは，ウォークが PFN を見つけられなかった
  ときだけ必要である．
\end{keypoint}

\section{TLB の構成と性能}
\label{sec:l13-tlb-org}

TLB の構成は，その役割から導かれる．\source{Lesson 13，動画
25--26；H\&P §B.4．}
\begin{description}
  \item[連想度] TLB ミスは1回ごとに数回のメモリ読み出しからなるウォークを
    要するので，TLB はフルアソシアティブか，非常に連想度の高いセット
    アソシアティブとする．これだけ少ないエントリ数であれば比較器の費用は
    許容でき，変換どうしがエントリを奪い合うこともなくなる．
  \item[サイズ] TLB のヒット率はキャッシュと同程度であるべきなので，
    キャッシュが保持するデータのページを覆えなければならない
    （\cref{ex:l13-tlb-size}）．これには 64 個から 512 個ほどのエントリが
    必要である．すべてのアクセスで参照されるので，L1 キャッシュのヒットと
    同程度に高速でもなければならず，そのためこれ以上大きくできない．\margin{Intel の最適化マニュアルによれば Skylake は，4~KB ページ用の第1
    レベルデータ TLB が 64 エントリ，命令用が 128 エントリ，その後ろに
    1536 エントリ 12 ウェイの第2レベル TLB をもつ．}
  \item[階層] それで十分なメモリを覆えない場合は，キャッシュを階層化するのと
    同じように（第14章），数百から数千のエントリをもつ，より大きく低速な
    第2レベルの TLB を第1レベルの後ろに置く．これは第1レベルでミスしたときに
    だけ参照され，両方の TLB がミスしたときにだけウォークが必要になる．
\end{description}

\begin{example}
  \label{ex:l13-tlb-performance}
  \cref{ex:l13-access-time} のプロセッサ（どのメモリ読み出しも AMAT は 4
  サイクル，4段ページテーブル）に，1 サイクルで参照でき，アクセスの
  \SI{1}{\percent} でミスする TLB を与える．どのアクセスも探索とデータ
  アクセスの代償を払い，100 回に1回は $4 \times 4 = 16$ サイクルのウォークの
  代償も払うので，アクセス1回は平均して 20 サイクルではなく
  $1 + 0.01 \times 16 + 4 = 5.16$ サイクルかかる．変換の費用は 16 サイクル
  ではなく 1.16 サイクルとなり，その大半は探索の分である．第14章では，この
  探索をキャッシュの探索と重ね合わせる．\margin{この式は \cref{eq:l12-amat} と
  同じ形をしている．探索時間 $+$ ミス率 $\times$ ペナルティであり，ウォークが
  ペナルティにあたる．TLB はキャッシュであり，同じように評価する．}
\end{example}

\begin{keypoint}[章のまとめ]
  仮想記憶は，各プログラムの仮想アドレス空間のページを物理メモリのフレームに
  対応づけ，残りをディスクに置く．フレームのないページへのアクセスは
  ページフォールトである．変換は，仮想ページ番号を，ページテーブルから得た
  フレーム番号で置き換え，オフセットはそのまま残す
  （\cref{eq:l13-translation}）．多段ページテーブルは，単一レベルのテーブルが
  使われないアドレス空間に費やすエントリを不要にする
  （\cref{eq:l13-flat-size}）．ページサイズはテーブルの大きさと内部断片化の
  トレードオフを決める．どの段もアクセス1回につきメモリ読み出しを1回加えるので，
  小さく連想度の高い TLB が変換をキャッシュし，ページテーブルウォークは TLB
  ミスのときだけ行われる．
\end{keypoint}

\end{document}
